
Devátý případ užití je specifikován v tabulce č. \ref{tab:usecase:UC09}.

\begin{UseCaseTable}
  {Správa životního cyklu senzoru}
  {UC09}
  {F09}
  {Změna provozního stavu senzoru nebo jeho přemístění do jiné místnosti}
  {Provozovatel monitorovaného prostoru}
  {}
  {Uživatel je přihlášen, senzor je v systému evidován}
  {Senzor má nový stav nebo je přiřazen k jiné místnosti}

\UseCaseScenario{Základní scénář -- změna stavu senzoru}
\UseCaseStep{uživatel}{V detailu evidovaného senzoru zvolí akci pro změnu provozního stavu.}
\UseCaseStep{systém}{Nabídne dostupné stavy (aktivní, pozastavený, vyřazený).}
\UseCaseStep{uživatel}{Zvolí nový stav a volbu potvrdí.}
\UseCaseStep{systém}{Aktualizuje stav senzoru a zobrazí potvrzení.}

\UseCaseScenario{Alternativní scénář -- přemístění senzoru}
\UseCaseStepCustom{A1}{uživatel}{V detailu evidovaného senzoru zvolí akci pro přemístění do jiné místnosti.}
\UseCaseStepCustom{A2}{systém}{Vrátí formulář s výběrem cílové místnosti z místností evidovaných uživatelem.}
\UseCaseStepCustom{A3}{uživatel}{Zvolí cílovou místnost a volbu potvrdí.}
\UseCaseStepCustom{A4}{systém}{Aktualizuje přiřazení senzoru k místnosti a zobrazí potvrzení.}

\end{UseCaseTable}
